%% Example CV - Vasudha Sathyapriyan
%% Comprehensive overview of all competencies, experience, and achievements.
%% Use as a master reference when tailoring role-specific CVs.
%%
%% SETUP: Run /setup to populate this with your actual information.
%% Compile with: lualatex main_example.tex
%% (pdflatex often fails on modern MiKTeX with fontawesome5 font-expansion errors.)

\documentclass[11pt,a4paper,sans]{moderncv}
\moderncvstyle{banking}
\moderncvcolor{blue}

% Force both first and last name AND section headings to render in moderncv
% blue (color1). Default banking on lualatex+MiKTeX leaves these black, which
% looks inconsistent with the rest of the blue accent scheme.
\renewcommand*{\firstnamestyle}[1]{{\fontsize{34}{36}\bfseries\upshape\color{color1}#1}}
\renewcommand*{\lastnamestyle}[1]{{\fontsize{34}{36}\bfseries\upshape\color{color1}#1}}
\renewcommand*{\sectionstyle}[1]{{\sectionfont\color{color1}#1}}

\usepackage[utf8]{inputenc}
\usepackage{hyperref}
\hypersetup{
    colorlinks=true,
    linkcolor=blue,
    filecolor=magenta,
    urlcolor=blue,
    pdftitle={Vasudha Sathyapriyan - CV},
    pdfpagemode=FullScreen,
}
\usepackage[scale=0.80]{geometry}
\usepackage{import}

% personal data
\name{Vasudha}{Sathyapriyan}
\address{Copenhagen, Denmark}{}{}
\phone[mobile]{+45-52696102}
\email{vasudhasathya@gmail.com}
\extrainfo{\href{https://www.linkedin.com/in/vasudha-sathya}{LinkedIn}, \href{https://vsathyapriyan.github.io/}{Portfolio}}

\begin{document}

\makecvtitle

% ============================================================
%     PROFILE STATEMENT
% ============================================================

\vspace{6pt}
\small{Signal processing engineer completing a PhD in speech enhancement, working across classical DSP and deep learning: adaptive beamforming, multi-microphone noise reduction, and DNN-based mask estimation. Experienced taking algorithms from research to deployed, real-time products using MATLAB, Python and PyTorch. 6 peer-reviewed publications and 2 patents.}

% ============================================================
%     CORE COMPETENCIES
% ============================================================

\section{Core Competencies}
\vspace{1pt}
\begin{itemize}

\item \textbf{Signal Processing \& Estimation}: Statistical signal processing, adaptive beamforming, multi-microphone noise reduction, detection and estimation theory, TDOA estimation

\item \textbf{Constrained Optimization}: Formulating noise-reduction and cue-preservation problems as constrained optimization, solved via Lagrangian methods

\item \textbf{Machine Learning for Audio}: Training and evaluating DNNs for mask estimation and post-filtering (PyTorch), benchmarked against classical and state-of-the-art methods

\item \textbf{Algorithm-to-Product Development}: Taking algorithms from research and derivation through acoustic measurement, simulation, real-time testing, and clinical validation to shipped hearing instruments

\item \textbf{Software Development}: Python, MATLAB, C++, Git, Linux, Agile methodology, system design

\end{itemize}

% ============================================================
%     PROFESSIONAL EXPERIENCE
% ============================================================

\section{Professional Experience}
\vspace{3pt}
\begin{itemize}

% --- Most Recent Role ---
\item{\cventry{2024--Present}{DSP Engineer}{Demant A/S (Oticon)}{Smørum, Denmark}{}{\vspace{1pt}
\begin{itemize}
    \item Develop and optimize DSP algorithms for hearing instruments across the full product lifecycle, from concept to commercial launch
    \item Research and implement DSP algorithms for adaptive beamforming and single- and multi-microphone noise reduction
    \item Develop data-driven machine learning methods for speech enhancement: train and evaluate DNNs for mask estimation in single-channel post-filtering (PyTorch), benchmarked against classical and state-of-the-art methods
    \item Verify algorithm performance through acoustic measurements, simulations, real-time testing, and clinical data analysis; define system requirements and test cases for deployment
\end{itemize}}}

\vspace{3pt}

% --- Previous Role ---
\item{\cventry{2019}{Graduate Research Intern}{Bosch Security and Safety Systems}{Eindhoven, Netherlands}{}{\vspace{1pt}
\begin{itemize}
    \item Investigated and improved beamforming algorithms for microphone arrays in conference room systems, validating the proposed algorithm against system requirements
\end{itemize}}}

\vspace{3pt}

% --- Earlier Role ---
\item{\cventry{2017--2018}{Software Engineer}{Visteon Corporation}{Chennai, India}{}{\vspace{1pt}
\begin{itemize}
    \item Developed embedded software for automotive instrument clusters, adhering to production coding standards
    \item Delivered products using Agile methodology, working with cross-functional domain heads to turn requirements into design and implementation
\end{itemize}}}

\end{itemize}

% ============================================================
%     EDUCATION
% ============================================================

\section{Education}
\vspace{1pt}
\begin{itemize}

\item{\cventry{2021--Present}{PhD in Electrical and Electronics Engineering (Industrial)}{Aalborg University / Demant A/S (Oticon)}{Smørum, Denmark}{}{\vspace{1pt}
Thesis: ``Speech enhancement in hearing aid applications using remote microphones.'' Enrolled at Aalborg University (Marie Skłodowska-Curie Fellowship), based full-time at Demant A/S (Oticon). Thesis in final review, defense expected October 2026. Visiting Ph.D. Student at Aalborg University (Oct. 2023 -- Feb. 2024) and Imperial College London (Oct. 2022 -- Feb. 2023).
}}

\vspace{3pt}

\item{\cventry{2018--2020}{MSc in Electrical and Electronics Engineering (Cum Laude)}{Delft University of Technology}{Delft, Netherlands}{}{\vspace{1pt}
Master's thesis on noise reduction in hearing aids with binaural cue (ILD/ITD) preservation.
}}

\vspace{3pt}

\item{\cventry{2013--2017}{B.E. in Electrical and Electronics Engineering}{SSN College of Engineering}{Chennai, India}{}{}}

\end{itemize}

% ============================================================
%     LANGUAGES
% ============================================================

\section{Languages}
\vspace{1pt}
\begin{itemize}
\item English (native), Tamil (native), Kannada (native), Hindi (native).
\end{itemize}

% ============================================================
%     PUBLICATIONS (optional)
% ============================================================

\section{Publications}
\vspace{1pt}
\begin{itemize}
\item V. Sathyapriyan et al. (2025). Head-Steered Channel Selection for Hearing Aid Applications Using Remote Microphones. IEEE Access, vol. 13.
\item V. Sathyapriyan et al. (2025). Binary Estimator Selection Methods for Hearing Aids With a Remote Microphone. IEEE Access, vol. 13.
\item V. Sathyapriyan et al. (2024). Speech Enhancement in Hearing Aids Using Target Speech Presence Estimation Based on a Delayed Remote Microphone Signal. IEEE ICASSP.
\item V. Sathyapriyan et al. (2023). Speech Enhancement using Binary Estimator Selection Applied to Hearing Aids with a Remote Microphone. IEEE ICFSP.
\item V. Sathyapriyan et al. (2022). A Linear MMSE Filter Using Delayed Remote Microphone Signals for Speech Enhancement in Hearing Aid Applications. IEEE IWAENC.
\item V. Sathyapriyan, M. Çaliş, R. C. Hendriks (2021). Binaural beam-forming with dominant spatial cue preservation for hearing aids. EUSIPCO.
\end{itemize}

% ============================================================
%     PATENTS
% ============================================================

\section{Patents}
\vspace{1pt}
\begin{itemize}
\item V. Sathyapriyan et al. A hearing aid system for determining an audio signal of interest. European Patent Appl. EP 24185393.6 (2024).
\item M. S. Pedersen, J. Jensen, N. R. Rohde, V. Sathyapriyan. Hearing system comprising a hearing aid and an external processing device. EP 4250765 A1 (2023).
\end{itemize}

% ============================================================
%     REFERENCES
% ============================================================

\section{References}
\vspace{1pt}
\begin{itemize}
\item{Available upon request.}
\end{itemize}

\end{document}
